\section{Introduction}
\label{sec:introduction}

\paragraph{Real-world privacy relevance.}
This paper addresses the right-to-erasure obligation of GDPR Article~17 as applied to
federated models under EDPB Opinion 28/2024 on AI systems \cite{edpb2024}, which
requires data controllers to evidence that a departing client's residual influence on
any downstream model is negligible. In cross-institutional and cross-device federated
learning deployments (medical, financial, keyboard, smart-home), this obligation is
operationally difficult: the server cannot observe raw client data by design, and
retraining from scratch is prohibitive. The protocol proposed here, FORGE, provides
a formal $(\varepsilon, \delta)$-indistinguishability guarantee that a supervisory
authority can verify without the controller publishing a pre-erasure model, and
without requiring the departing client to remain online. The guarantee, the
verification protocol (Section~\ref{sec:method-erasure}), and the cross-release
composition analysis (Section~\ref{sec:theory}) are the components that make the
protocol deployable as a compliance mechanism for live FL systems rather than as a
theoretical construction.

The right to erasure obliges a data controller to remove personal data on request and,
under recent guidance, to evidence that residual influence on any downstream model is
negligible \cite{edpb2024}. In Federated Learning (FL), and especially in Federated
Continual Learning (FCL), operationalising this obligation is complicated by the very
properties that make FL privacy-preserving in the first place. Secure aggregation
prevents the server from observing individual client updates \cite{bonawitz2017};
asynchronous or streaming continual learning spreads a client's influence across weights,
aggregated statistics, replay buffers, and distilled knowledge \cite{yoon2021}; and
empirical work has shown that secure aggregation is not, on its own, sufficient
protection against membership inference \cite{pillutla2023}. Any credible Federated
Unlearning (FU) procedure must therefore not merely alter the model, but do so in a way
that a supervisory authority, as well as the departing client themselves, can verify
without introducing new avenues of leakage.

Existing FU methods trade among three costs, and no prior work eliminates all three
simultaneously. First, \emph{history-based} methods reconstruct or calibrate a model
without the departing client using stored per-round updates or gradients
\cite{liu2021federaser,fraboni2024sifu,zhang2023fedrecovery}. This gives the tightest
fidelity to a retrained model, but the storage itself is at odds with secure aggregation
and, under EDPB Opinion 28/2024, may not qualify as anonymous processing. In
cross-device deployments with millions of clients, it is also often infeasible. Second,
\emph{interactive} methods avoid storage by having the departing client re-engage with
the server: gradient ascent on their local data with a distance constraint from the
global model \cite{halimi2022}. This is elegant but assumes the client's continued
participation, which is an assumption that fails precisely when the client is exercising
a right to leave. Third, \emph{heuristic} methods (weight negation, pruning,
random-label perturbation) sidestep both storage and cooperation but offer no formal
guarantee a controller can evidence \cite{khalil2025not,gu2024fedau}. Recent
memorisation-scoped approaches \cite{wei2026fedmemprune} give strong empirical fidelity
but do not currently ship with a provable indistinguishability bound.

A separate, less-recognised failure mode is that the unlearning procedure itself leaks.
Publishing (or letting a client observe) a before/after model pair enables membership
inference attacks (MIA) that outperform attacks on the original model
\cite{chen2021unlearningleaks}, and dedicated federated attacks such as the Federated
Unlearning Inversion Attack (FUIA) exploit exactly this signal \cite{fuia2025}. Any
method that requires a comparable ``after'' model to be released must be evaluated
against this channel.

In this paper, we propose FORGE (\textbf{For}ward-forgetting via decayed aggregation and
pre-committed erasure summaries), an FCL protocol under which the erasure guarantee is a
property of the training dynamics, rather than a procedure invoked at request time. The
protocol has two ingredients. (1)~\emph{Bounded-influence decay.} We formalise a
training-time property called $(\varepsilon, \delta, T)$-forward forgetting: a client
that departs at round $r$ is $(\varepsilon, \delta)$-indistinguishable from a
counterfactual training run that never included them, by round $r + T$. Under a decayed
aggregation protocol with per-round clipping constant $C$ and calibrated server noise
$\sigma$, we exhibit an explicit rate $\varepsilon_1(T)$. This is the FL-native analogue
of the DP-as-unlearning line \cite{guo2020certified,sekhari2021remember,neel2021descent},
but time-indexed rather than budgeted once-for-all. (2)~\emph{Pre-committed differentially
private (DP) summary channel.} During training, each client submits alongside their
gradient a small DP summary $\sigma_i$ of their local data distribution. At erasure time,
the server uses $\sigma_{\forget}$ to perform a non-interactive subtraction of the
client's memorised residual, at additional cost $\varepsilon_2$. No per-round updates
are stored, and no interaction with the client is required at or after the request.
Because the feature-block aggregation and the summary channel use independent Gaussian
noise sources, and because the summary-subtraction step at erasure time operates only
on the head, adaptive Gaussian-DP composition applies and gives a
$\sqrt{\mu_1(T)^2 + \mu_\summary^2}$-GDP joint release as the overall formal guarantee.

\paragraph{Contributions.}
\begin{enumerate}[leftmargin=*,itemsep=2pt]
\item A new privacy definition: $(\varepsilon, \delta, T)$-forward forgetting, the first
per-client, time-indexed unlearning guarantee for federated (continual) learning
(Def.~\ref{def:forward-forgetting}).
\item A decayed aggregation FCL protocol with per-round clipping and calibrated server
noise that achieves forward forgetting at an explicit rate
(Alg.~\ref{alg:decayed-aggregation}).
\item A pre-committed DP summary construction that enables non-interactive server-side
subtraction of a departed client's memorised residual, with no per-round update storage
and no request-time client cooperation (Alg.~\ref{alg:erasure}).
\item A composition theorem (Thm.~\ref{thm:composition}) that yields a
$\sqrt{\mu_1(T)^2 + \mu_\summary^2}$-GDP guarantee by adaptive Gaussian-DP
composition, enabled by the independent Gaussian noise sources of the feature-block
and summary mechanisms in Algorithm~\ref{alg:decayed-aggregation}.
\item An empirical evaluation whose primary benchmark is CIFAR-10 non-IID
stationary FL (ResNet-18, 100 clients, 200 rounds, Dirichlet $\alpha{=}0.3$, three
seeds) at realistic deployment utility, with a class-incremental stress test
comparing FORGE against eleven baselines including an $\varepsilon$-sweep of
FedRecovery across $\varepsilon \in \{1, 2, 5, 10, \infty\}$, plus cross-dataset
confirmation on Split-CIFAR-100 and FEMNIST. We evaluate with Yeom and
shared-shadow Likelihood Ratio Attack (LiRA) membership-inference attacks as well
as the FUIA gradient-inversion attack; the full benchmark is reproducible from a
single command on consumer M-series hardware.
\end{enumerate}

Table~\ref{tab:positioning} summarises how FORGE maps onto the three FU pain points
relative to representative prior work. To the best of our knowledge, FORGE is the first
FU protocol to resolve all three simultaneously.

\begin{table}[t]
\centering
\small
\setlength{\tabcolsep}{4pt}
\begin{tabular}{lccc}
\toprule
Method & No hist. & No coop. & Formal guar. \\
\midrule
Retrain from scratch & \checkmark & \checkmark & exact \\
FedEraser \cite{liu2021federaser} & $\times$ & \checkmark & $\times$ \\
SIFU \cite{fraboni2024sifu} & $\times$ & \checkmark & \checkmark \\
FedRecovery \cite{zhang2023fedrecovery} & $\times$ & \checkmark & \checkmark (DP) \\
Halimi et al.\ \cite{halimi2022} & \checkmark & $\times$ & $\times$ \\
FedMemPrune \cite{wei2026fedmemprune} & \checkmark & \checkmark & $\times$ \\
\textbf{FORGE (ours)} & \checkmark & \checkmark & \checkmark (DP) \\
\bottomrule
\end{tabular}
\caption{Positioning of FU methods against the three cost axes: ``No hist.'' = no per-round history storage; ``No coop.'' = no client cooperation at request time; ``Formal guar.'' = explicit $(\varepsilon,\delta)$ bound to the retrain counterfactual.}
\label{tab:positioning}
\end{table}

\paragraph{Paper organization.}
The rest of the paper is structured as follows. Section~\ref{sec:related} discusses
related work. Section~\ref{sec:prelim} introduces the system model, the threat model,
the forward-forgetting definition, and the problem statement. Section~\ref{sec:method}
describes the FORGE protocol and its two algorithms. Section~\ref{sec:theory} states and
proves the three main theoretical results. Sections~\ref{sec:experiments},
\ref{sec:results}, and \ref{sec:results-femnist} describe the experimental setup and
report results on Split-CIFAR-10 and FEMNIST, respectively. Section~\ref{sec:discussion}
interprets the main findings. Section~\ref{sec:limitations} summarises the scope of the
evaluation. Finally, Section~\ref{sec:conclusion} concludes the paper and presents
future research directions.
